\chapter{本体论基础理论}
\chapterart{ch02}{类层次（TBox）与实例（ABox）的两层世界}

\begin{introbox}
\textbf{【导读】}
如果把本体比作一座建筑，本章就是它的承重结构：类、属性、公理、实例。
这四样东西看似简单，组合起来却能表达一个领域的完整知识骨架。
本章先把四块积木逐一拆开细看——类层次怎么搭、属性特性怎么选、公理怎么写、实例怎么管，
再引入两套「数学语言」——描述逻辑与一阶逻辑，说明 OWL 为什么选择前者作为地基。
最后讨论两个容易被忽视、却直接决定系统行为的问题：
没有记录的事实算不算「假」（开放世界与封闭世界）？
新知识到来时，旧结论会不会被推翻（单调与非单调）？
读完本章，你应当能看懂任何一个 OWL 本体的骨架，
并且对「这条公理会让推理器做出什么」建立起可靠的直觉。
\end{introbox}

\chapterbinding{semantica.chapter\_packages.vol1.ch02}{partial}{blocked}
{CQ、OWA/CWA 查询、输入数据与受限正向规则场景已绑定；
\texttt{OE-V1-CH02-SCN-REASONING-MODES-001} 有精确 oracle。}
{当前只承诺正向、单调规则适配与 RDF/SPARQL 复算；完整 DL/tableau、默认逻辑、
一般非单调推理与本章 SHACL 均不支持，必须 fail closed。}

\section{四块积木总览：类、属性、公理、实例}

假设你刚接到一个任务：把一家机加工厂关于设备的知识「装进」信息系统，
让排产程序不再依赖老师傅的口头提醒。摆在你面前的是一份两百行的设备台账，第一行写着：
「EQ-2023-0117，3 号车床，数控车床，主电机功率 15.5 kW，位于 A 区 2 号工作单元」。
这一行朴素的记录里，其实已经藏着本体的全部四种构件。把它拆开，就是本章的起点。

\textbf{类}（Class）是概念的抽象。「数控车床」不是指这台具体的机器，
而是指一类机器——所有满足「数控车床」资格判据的设备。
台账里的「设备」「数控车床」，乃至没写出来的「加工设备」「材料」「工序」，都是类。
类回答的问题是：\textbf{它是什么}。

\textbf{属性}（Property）是连接的纽带。「主电机功率 15.5 kW」与「位于 A 区 2 号工作单元」
是两种性质不同的连接：前者把个体连到一个字面值（数字 15.5），
叫\textbf{数据属性}（Data Property，如 \texttt{hasPower}）；
后者把个体连到另一个个体——那个工作单元本身也是一个可以被继续描述的对象，
叫\textbf{对象属性}（Object Property，如 \texttt{locatedIn}）。
属性回答的问题是：\textbf{它与谁有关、有什么特征}。

\textbf{公理}（Axiom）是不允许违反的规则。它们在台账上看不见，却在每个人脑子里：
「车床是设备的一种」「一台设备同一时刻只在一个工作单元」「设备不是材料」。
把这些规则显式地写进本体，机器才第一次有了替你检查错误、推导新知识的依据。
公理回答的问题是：\textbf{什么是必然的、什么是不可能的}。

\textbf{实例}（Individual，也译个体）是落地的具体对象——编号 EQ-2023-0117 的那台 3 号车床，
世界上独一无二的存在。类描述「车床应当是什么样」，实例是「这一台车床」。
实例回答的问题是：\textbf{现实中到底有谁}。

四块积木中，类层次是最直观的部分。看 ch02 包内教学资产中的设备分类树：

\begin{codebox}
\codeline{\textcolor{cmtgray}{\#\ ==============================}}
\codeline{\textcolor{cmtgray}{\#\ 1.\ 类层次结构\ (Class\ Hierarchy)}}
\codeline{\textcolor{cmtgray}{\#\ ==============================}}
\codeline{\textcolor{cmtgray}{\#\ 类是本体的基础单元，代表领域内的一类事物}}
\codeline{\textcolor{cmtgray}{\#\ 子类会继承父类的全部属性与关系}}
\codeblank{}
\codeline{\textcolor{cmtgray}{\#\ 设备分类层次：}}
\codeline{\textcolor{cmtgray}{\#\ 设备\ (Equipment)}}
\codeline{\textcolor{cmtgray}{\#\ \ \ ├──\ 加工设备\ (ProcessingEquipment)}}
\codeline{\textcolor{cmtgray}{\#\ \ \ │\ \ \ ├──\ 数控机床\ (CNCMachine)}}
\codeline{\textcolor{cmtgray}{\#\ \ \ │\ \ \ │\ \ \ ├──\ 数控车床\ (CNCLathe)}}
\codeline{\textcolor{cmtgray}{\#\ \ \ │\ \ \ │\ \ \ └──\ 数控铣床\ (CNCMillingMachine)}}
\codeline{\textcolor{cmtgray}{\#\ \ \ │\ \ \ └──\ 手动机床\ (ManualMachine)}}
\codeline{\textcolor{cmtgray}{\#\ \ \ ├──\ 测量设备\ (MeasurementEquipment)}}
\codeline{\textcolor{cmtgray}{\#\ \ \ ├──\ 辅助设备\ (AuxiliaryEquipment)}}
\codeline{\textcolor{cmtgray}{\#\ \ \ └──\ 工装设备\ (ToolingEquipment)}}
\codeblank{}
\codeline{\textcolor{cmtgray}{\#\ 设备(Equipment)\ \(\rightarrow\)\ 加工设备\ /\ 测量设备\ /\ 辅助设备\ /\ 工装设备}}
\codeline{\textcolor{cmtgray}{\#\ \(\rightarrow\)\ 数控机床\ /\ 手动机床等细分}}
\codeblank{}
\end{codebox}

\assetsource{ch02}{core-concepts}

这棵树要竖着读：每一台数控车床都是数控机床，每一台数控机床都是加工设备，
每一台加工设备都是设备——沿树上行的每一步都是一次「是一种」（is-a）判断。
子类继承父类的全部属性与关系：为「加工设备」声明的「需要定期维护」，
自动适用于每一台数控车床，无须逐类重复。这种「说一次，处处生效」的经济性，
是层次结构的第一重价值；第二重价值要到下一节才能看清——层次让推理器能够替你\textbf{分类}。

拆解完毕，可以把台账第一行完整地翻译成本体语言了：

\begin{noteblock}
\textbf{台账一行的本体化翻译。}
类层面（设计一次）：\texttt{CNCLathe \(\sqsubseteq\) CNCMachine}、
\texttt{CNCMachine \(\sqsubseteq\) ProcessingEquipment}、
「设备与材料不相交」等公理。\par
实例层面（逐行生成）：\texttt{CNCLathe(EQ-2023-0117)}——类断言；
\texttt{hasPower(EQ-2023-0117, 15.5)}——数据属性断言；
\texttt{locatedIn(EQ-2023-0117, CellA2)}——对象属性断言，
其中 \texttt{CellA2} 是「A 区 2 号工作单元」这个个体。\par
台账其余一百九十九行，只需重复实例层面的翻译——类与公理一次设计，处处复用。
\end{noteblock}

习惯上把类与公理的集合叫 \textbf{TBox}（术语层，Terminology），
把实例断言的集合叫 \textbf{ABox}（断言层，Assertion）。
章首插图画的正是这两层世界：上层是稳定的概念树，由本体工程师精心设计，
演化以月为单位；下层是海量的实例点，由业务系统源源不断地产生，变化以秒为单位。
两层分离，知识的「骨架」与「血肉」就可以独立演化——这是本体区别于普通数据库表设计的
关键架构决策：改一张数据库表往往牵动所有应用，而 TBox 的一次精心修订，
ABox 里的千万条断言原地不动。第 7 章的「知识图谱」，本质上就是 TBox 加上大规模 ABox。

最后强调四块积木的整体性：没有属性的类只是一片标签云，
没有公理的类与属性只是随手连线的示意图，没有实例的本体则是永远不接触现实的空中楼阁。
一句话概括本章的对象：\textbf{本体 = 词汇（类与属性）+ 约束（公理）+ 事实（实例断言）}。

把四块积木对照到熟悉的工具上，有助于消化，也能暴露差异：
Excel 的一列像数据属性，数据库的外键像对象属性，表上的检查约束像公理，
一行记录像实例——但每一处「像」都藏着本质不同。
检查约束只在写入时拦截一次，公理随时参与推理、还能推出从未录入过的新知识；
外键只保证「指得到」，对象属性还能声明传递、对称等逻辑性质（2.3 节）；
数据库模式为某一个应用的存取效率而设计，本体为整个组织的共识与推理而设计。
所以同一个领域往往同时存在多个数据模型（订单系统一个、排产系统一个），
却应当只有一套本体——前者是各自的存储视角，后者是共同的语义契约。
接下来四节依次把四块积木放大细看，之后再回答「这一切的数学地基是什么」。

\section{类与类层次：把「是一种」说严格}

「数控车床是一种设备」，没有人会说错。但下面三个问题，几乎每个建模评审会都会吵起来：
「备用设备」是不是设备的子类？「主轴」是不是车床的子类？
「3 号车床」是不是数控车床的子类？——答案是三个「不是」。
直觉在简单场景里够用，在边界场景里失灵。本节把「是一种」说严格，给出可操作的判断工具。

\subsection{外延与内涵：类的两副面孔}

一个类有两副面孔。\textbf{外延}（extension）是它当下全部成员的集合——
比如全厂现役的 17 台数控车床；\textbf{内涵}（intension）是成员资格的判据——
「由数控系统控制、主轴带动工件旋转完成切削加工的机床」。
外延随时间变化：明天报废一台，就只剩 16 台；内涵相对稳定，可能十年不改一个字。

写进本体的是内涵——类的定义与公理；外延不直接维护，
由 ABox 断言加推理在运行时给出。初学者最常见的心智偏差，
是把类当成「成员清单」来管理——那是数据库视图干的事。
类回答的是「凭什么算」，不是「现在有谁」。

两副面孔甚至可能短暂重合而依然不同。
「当前位于 A 区的设备」与「编号以 A 开头的设备」，在某个时刻外延可能恰好相同
（A 区设备恰好都按区编号），但它们是两个类：判据不同、变化规律不同、该挂的公理也不同——
前一个跟着调度走，后一个跟着编码规则走。
推理器关心的、公理约束的，从来都是内涵。
分清这两副面孔，后面讨论的许多问题（定义类、自动分类、全称陷阱）才有立足点。

\subsection{subsumption 的集合语义}

子类关系的术语叫 \textbf{subsumption}（归类关系，记号 \(\sqsubseteq\)）。
\texttt{CNCLathe \(\sqsubseteq\) Equipment} 的精确含义是一个集合论命题：
\textbf{在任何可能的状态下}，数控车床的外延都是设备外延的子集——
不管工厂如何扩建、台账如何更新，凡是数控车床必是设备，没有例外。
于是检验一条子类边该不该画，有一句口诀：
「每一个 X，在任何情况下，都必然是 Y 吗？」只要能想出一个反例场景，这条边就不成立。

由此立刻得到一个重要推论：\textbf{继承不是数据拷贝，而是逻辑推论}。
「数控车床需要定期维护」不是系统在建类时把维护字段「复制」了一份，
而是推理器在被问到时现场推出来的必然结论：
凡数控车床是加工设备，凡加工设备需维护，故凡数控车床需维护。
推论无所谓「同步失效」，也永远不会和来源不一致。

也因此，本体的 subsumption 与面向对象编程的继承貌合神离，值得专门辨析。
OO 的继承是\textbf{实现复用}机制，子类可以覆写（override）父类行为；
本体的 subsumption 是\textbf{逻辑蕴含}，子类无法「取消」父类的任何约束——
如果你发现某个「子类」需要豁免父类的规则，那它根本就不是子类。
这一条判据可以在评审会上消灭一半的错误层次：
「免检产品」想豁免「产品必须检验」？那「免检产品」就不该挂在那个「产品」之下，
而是上层的「产品必须检验」这条公理本身写得过强。

\subsection{单继承还是多重继承}

逻辑上，一个类可以有多个父类：车铣复合加工中心既是车削类设备又是铣削类设备；
带在机测量功能的机床既是加工设备，又承担测量设备的职能。
描述逻辑对多重继承毫无障碍——\(\sqsubseteq\) 只是集合包含，
一个集合当然可以同时是多个集合的子集；OO 语言里令人头痛的「菱形继承」问题，
在逻辑语义下根本不存在（没有「方法冲突」，只有约束的并集）。

问题出在人身上。多重继承的层次图不再是一棵树，而是一张格状网：
主干不清晰，评审时没人能一眼看出「这个类到底挂在哪」；
修改一个上层类时，影响面顺着多条父链扩散，难以预估；
新人接手时，理解成本成倍增加。规模越大，这些成本越是主导因素。

工程上的成熟做法是折中：\textbf{主干保持单继承}——每个类指定一个「主」父类，
整张图打印出来仍然像一棵树，人能读；\textbf{交叉维度交给定义类}——
「高端设备」「移动设备」「待维护设备」这类切面式分类，
用充要条件定义（2.4 节的 \(\equiv\) 公理），让推理器把满足条件的类自动「归并」进来。
一句话：\textbf{人手画的层次保持简单，复杂的交叉分类让机器算}。
第 4 章讲 OWL 的定义类（defined class）语法时，会看到这一策略的落地写法。

\subsection{命名规范}

命名看似小事，却是本体可维护性的第一道防线。业界通行的惯例可以归纳为五条。
其一，类名用首字母大写的驼峰式\textbf{单数}名词：\texttt{CNCLathe}，
而不是 \texttt{cnc\_lathes}——类指「这一种」，不指「这一堆」。
其二，属性名用小写开头的动词短语：\texttt{producedBy}、\texttt{hasPower}、\texttt{locatedIn}，
读起来就是一句主谓宾。
其三，不把层次信息编码进名字：写 \texttt{CNCLathe}，
不写 \texttt{Equipment\_Processing\_CNC\_Lathe}——层次属于公理，
调整层次不应该牵连改名。
其四，机器标识与人类语言分离：IRI 里用稳定的英文或编号，
中文「数控车床」放在标签（\texttt{rdfs:label}）里；需要多语言支持，就挂多个标签。
其五，标识符一经发布不再更改：改的是标签，不是 IRI——
引用它的外部系统不计其数，IRI 就是契约。
顺带记一个高频反模式：把版本号写进类名（\texttt{EquipmentV2}）。
版本属于本体文件的元数据（由注释信息承载，见 2.3 节），不属于概念本身；
概念的语义真变了，就该是一个新概念、配新名字与新定义，而不是旧名字加后缀。

\subsection{常见层次设计错误}

回到本节开头的三个争议。它们各自代表一类高频错误，连同另外两类一起列成下表：

\begin{center}
\begin{tabularx}{\linewidth}{@{}lXX@{}}
\toprule
\textbf{错误模式} & \textbf{典型例子} & \textbf{病因与修正} \\
\midrule
实例当子类 & 「3 号车床」建成 \texttt{CNCLathe} 的子类 &
它没有自己的成员资格判据，它本身就是一个成员。修正：类断言
\texttt{CNCLathe(Lathe\_001)}。判别法：问「它还能有实例吗」 \\
部分当子类 & 「主轴」建成「车床」的子类 &
主轴不「是一种」车床，而是车床的组成部分；is-a 与 part-of 混淆。
修正：对象属性 \texttt{partOf}。判别法：「每一根主轴都是一台车床吗」 \\
角色当类 & 「备用设备」「在修设备」建成 \texttt{Equipment} 的子类 &
「备用」是设备暂时扮演的角色，随调度一天三变；
类成员资格不该如此易变。修正：状态属性 \texttt{hasStatus}，或按需用定义类 \\
组合爆炸 & 「高端数控车床」「低端数控车床」「高端数控铣床」…… &
把每个维度组合都建成类，类数随维度数指数增长。
修正：维度做属性，交叉概念用定义类按需生成 \\
维度混层 & 「加工设备」（按功能分）与「移动设备」（按移动性分）并列为兄弟 &
同一层混用两个划分维度，成员归属模糊、兄弟无法互斥。
修正：每一层只用一个划分维度，维度本身显式记录 \\
\bottomrule
\end{tabularx}
\end{center}

这些错误有一个共同特征：在几十个类的小本体里毫无症状，
等规模上来、推理器开动之后，分类结果开始「说胡话」，而那时层次已经被几十个应用引用。
所以层次评审值得在早期投入真功夫。第 3 章将介绍的 \textbf{OntoClean} 方法，
正是把本节这些靠直觉的判断（「备用」是角色不是类、「主轴」是部分不是种类）
升级为一套可执行的元属性检查表——刚性、同一性、统一性——让层次评审从各执一词变成按表走查。

\subsection{层次的形状：深度、宽度与电梯测验}

除了对错，层次还有「体形」问题。宽度上，一个类的直接子类如果多到十几个，
通常说明缺了中间层级，或者同一层混进了多个划分维度（上表第五种错误）；
常见的经验做法是：兄弟数明显膨胀时审视是否需要分组——但分组必须有业务含义，
为凑数而造的中间类只会徒增导航成本。深度上，过深的链条（动辄八九层）
让每次建模讨论都要从头爬树，过浅（只有两三层）则往往说明抽象不足、公理无处安放。

一个朴素而有效的体检手段是「电梯测验」：任取一个实例，沿父类链一路上行，
每一步大声读出「……是一种……」——
「这台机器是一种数控车床，是一种数控机床，是一种加工设备，是一种设备」。
只要有一步读起来别扭，那一层就值得复查。
层次的形状没有铁律，但有一条硬标准：\textbf{每一层的划分维度必须单一且显式}——
回头看上表，这正是多数错误的总根源。

\section{属性的世界：两种属性与七种特性}

只有类层次的本体是一棵「死树」：它能回答「数控车床是不是设备」，
却回答不了调度员真正关心的问题——这台设备能加工什么材料？位于哪个工作单元？由谁操作？
让知识活起来的是属性。本节先分清两种属性，
再纠正一个关于域与值域的经典误解，最后把七种属性特性逐一过堂。

\subsection{对象属性与数据属性}

\textbf{对象属性}连接个体与个体，\textbf{数据属性}连接个体与字面值（literal）。
区分两者有一个简单的追问：\textbf{这个值本身还需要被描述吗？}
「A 区 2 号工作单元」有面积、有位置、有所属车间——它应当是个体，用对象属性连接；
「15.5」就只是一个数，数据属性足矣。先看数据属性在台账场景下的样子：

\begin{codebox}
\codeline{\textcolor{cmtgray}{\#\ ==============================}}
\codeline{\textcolor{cmtgray}{\#\ 2.\ 数据属性\ (Data\ Properties)}}
\codeline{\textcolor{cmtgray}{\#\ ==============================}}
\codeline{\textcolor{cmtgray}{\#\ 描述实体的具体属性值（字面量）}}
\codeblank{}
\codeline{Equipment.serialNumber\ \(\rightarrow\)\ 字符串}
\codeline{Equipment.weight\ \(\rightarrow\)\ 浮点数（单位：kg）}
\codeline{Equipment.installDate\ \(\rightarrow\)\ 日期}
\codeline{Product.price\ \(\rightarrow\)\ 浮点数（单位：元）}
\codeline{Process.duration\ \(\rightarrow\)\ 整数（单位：分钟）}
\codeblank{}
\end{codebox}
\color{black}

\assetsource{ch02}{core-concepts}

每条数据属性都标明了取值类型，实践中对应 XSD 数据类型
（\texttt{xsd:string}、\texttt{xsd:float}、\texttt{xsd:date} 等，第 4 章正式介绍）。
这份示例里还藏着两个工程细节。第一，\textbf{单位藏在注释里}（kg、元、分钟）——这是隐患：
注释不参与推理，两套系统一个用千克一个用吨，机器毫无察觉。
更稳妥的做法是把单位显式化：或者属性名自带单位（\texttt{weightKg}），
或者把「量值」建成带数值与单位两个属性的小对象（第 6 章的工业案例会见到这种模式）。
第二，\textbf{不要把本该是对象属性的信息塞进字符串}：
如果 \texttt{locatedIn} 的值是字符串「A区」，
那么「A 区在几号厂房」「A 区还有哪些设备」就永远查不下去了——字符串是知识图谱的死胡同。

再看对象属性：

\begin{codebox}
\codeline{\textcolor{cmtgray}{\#\ ==============================}}
\codeline{\textcolor{cmtgray}{\#\ 3.\ 对象属性\ (Object\ Properties)}}
\codeline{\textcolor{cmtgray}{\#\ ==============================}}
\codeline{\textcolor{cmtgray}{\#\ 描述实体之间的关系}}
\codeblank{}
\codeline{Product.producedBy\ \(\rightarrow\)\ Equipment（产品由哪个设备生产）}
\codeline{Equipment.locatedIn\ \(\rightarrow\)\ WorkCell（设备位于哪个工作单元）}
\codeline{Process.requires\ \(\rightarrow\)\ Material（工序需要哪些材料）}
\codeline{Process.precedes\ \(\rightarrow\)\ Process（工序之间的先后顺序）}
\codeline{Worker.operates\ \(\rightarrow\)\ Equipment（工人操作哪些设备）}
\codeblank{}
\end{codebox}

\assetsource{ch02}{core-concepts}

五条属性勾勒出制造领域的关系骨架：产品与设备、设备与空间、工序与物料、
工序与工序、人与设备。读对象属性清单的正确姿势，是把每一条读成一个\textbf{可问的问题方向}：
有了 \texttt{producedBy}，就能问「这批产品出自哪台设备」；
有了 \texttt{precedes}，就能问「这道工序之前还有哪些工序」。
对象属性还可以声明\textbf{逆属性}（inverse property）：
\texttt{producedBy} 的逆是 \texttt{produces}，断言任意一个方向，另一个方向自动成立，
查询时两头都顺手。

严格说来，OWL 里还有第三种属性：\textbf{注释属性}（Annotation Property）。
\texttt{rdfs:label}（显示名）、\texttt{rdfs:comment}（说明文字），
以及出处、版本、责任人等管理元数据，都靠它挂载。
注释属性不参与逻辑推理——推理器对注释一律视而不见——却深度参与\textbf{治理}：
2.2 节「IRI 与标签分离」的纪律、2.4 节将提到的「每条公理写明业务出处」，
落点都是注释属性。一个没有注释的本体，半年之后连作者自己都不敢改。

属性之间同样可以组成层次：\textbf{子属性}（subPropertyOf）声明一条属性蕴含另一条，
例如 \texttt{coldRepairs \(\sqsubseteq\) maintains}——凡大修记录都是维护记录。
断言时走具体的子属性，查询时按需选粗细：
问「谁维护过这台设备」，大修、小修、保养一网打尽；问「谁大修过」，又能精确命中。
本节稍后会看到属性层次与传递性配合的「双属性模式」，
2.6 节则把这类属性公理正式归入 RBox。

\subsection{域与值域：是分类，不是校验}

每个属性可以声明\textbf{域}（domain，主语应属的类）与\textbf{值域}（range，宾语应属的类）：
\texttt{producedBy} 的域是 \texttt{Product}，值域是 \texttt{Equipment}。
到这里为止一切都像数据库的外键约束——而这正是从传统 IT 转来的工程师
最容易踩、踩得最痛的一个坑，值得用一个完整例子讲透。

设想有人误把一张订单连了上去：\texttt{producedBy(Order\_77, Robot\_R2)}。
数据库思维的预期是：报错，拒绝写入，因为 \texttt{Order\_77} 不是 \texttt{Product}。
OWL 推理器的实际行为是：安静地接受这条断言，
然后\textbf{推出} \texttt{Product(Order\_77)}——
「既然它出现在 \texttt{producedBy} 的主语位置，那它一定是个产品」。
没有报错，没有警告，你的订单从此被分类为产品，并开始享受产品的一切推理待遇。

这就是必须记住的一句话：\textbf{域与值域是推理公理，不是校验规则}。
它们的语义是「凡出现在该属性主语位置的个体都属于域类，凡出现在宾语位置的都属于值域类」，
告诉推理器如何\textbf{归类}陌生个体，而不是如何\textbf{拒绝}可疑数据。
这一设计与开放世界假设（2.8 节）一脉相承：知识默认不完备，
一个出现在 \texttt{producedBy} 主语位置却没有声明类型的个体，
多半只是「还没来得及声明」，推理器替你补上分类，恰恰是在帮忙。

那么错误何时才会被抓住？只有当推出的分类与已有知识\textbf{矛盾}时：
若本体声明了订单与产品不相交（\texttt{Order \(\sqcap\) Product \(\sqsubseteq\) \(\bot\)}），
推出 \texttt{Product(Order\_77)} 的瞬间知识库即告不一致，红灯亮起。
这里再次看到下一节的主题：想让错误暴露，公理必须足够强。
至于真正意义上的「输入校验」，那是另一类需求，
应交给专门的约束语言（如 SHACL）或入库前的数据管道检查，而不要指望域与值域。

\subsection{属性特性的全谱系}

OWL 允许为属性声明一组逻辑特性。每声明一种，
等于和推理器签下一条合同：它会据此多推一些结论，也会据此多抓一类矛盾。
七种特性逐一过堂，每种配一个制造业例子和一个误用后果。

\textbf{函数性}（Functional）：每个主体\textbf{至多}一个值。
\texttt{locatedIn} 是典型——一台设备同一时刻只在一个工作单元；
数据属性同样可以函数化，如 \texttt{hasSerialNumber}。
注意「至多」不是「恰好」：函数性不保证值存在，想要「恰好一个」还需基数公理（2.4 节）。
误用后果出人意料：把一名工人可以操作多台设备的 \texttt{operates} 错误地声明为函数性，
当 ABox 出现 \texttt{operates(Wang, Lathe\_001)} 与 \texttt{operates(Wang, Mill\_002)} 时，
推理器\textbf{不报错}，而是推出 \texttt{Lathe\_001} 与 \texttt{Mill\_002} 是同一个体——
车床与铣床被合并成一台机器。为什么不报错？因为 OWL 不采用唯一名假设（2.5 节详述），
两个名字默认可以指同一个东西。\textbf{函数性遇到多值，结论是合并而非冲突}——
除非那两个个体已被显式声明互异。

\textbf{逆函数性}（InverseFunctional）：一个值\textbf{至多}属于一个主体，
即值反过来唯一确定主体。「一个 RFID 标签至多贴在一台设备上」：
\texttt{hasRFIDTag} 应声明逆函数。两个「设备个体」共享同一个标签，推理器推出它们是同一台——
这是数据融合的利器：ERP 与 MES 各自建了设备个体，
只要挂上同一个标签个体，自动合并，无须人工对账。
同一机制反过来就是风险：错误的共享值引发错误合并，且合并的连锁影响很难回滚。
还有一条标准细节：OWL 2 DL \textbf{不允许数据属性声明逆函数}——
想拿序列号字符串当「天然主键」走不通，正路是第 4 章将介绍的键公理（\texttt{HasKey}）。

\textbf{传递性}（Transitive）：由 R(a,b) 与 R(b,c) 推出 R(a,c)。
空间包含是教科书案例：设备位于工位、工位位于车间、车间位于厂区，
\texttt{locatedIn} 声明传递后，「这台设备在哪个厂区」一步可达，无须应用层循环展开；
\texttt{partOf} 同理。误用后果：把「直接前驱」\texttt{directlyPrecedes} 声明为传递，
「直接」二字的语义瞬间被摧毁——传递闭包让所有间接前驱都变成了「直接前驱」。
一般规律：\textbf{名字里带「直接」「相邻」的属性绝不能传递}。
工程上常用双属性模式：\texttt{directlyPrecedes} 负责记录原始数据（不传递），
声明它是传递属性 \texttt{precedes} 的子属性，让推理器在 \texttt{precedes} 上算闭包——
数据干净，推理也不缺。

\textbf{对称性}（Symmetric）：R(a,b) 成立则 R(b,a) 成立。
\texttt{adjacentTo}（工位相邻）、\texttt{connectedTo}（管路连通）是正例：
断言一个方向，另一个方向自动成立，维护量减半。
误用后果：把 \texttt{precedes} 声明对称——「焊接先于热处理」自动推出「热处理先于焊接」，
全部时序约束作废，排程逻辑崩塌。更隐蔽的误用是想当然：
「协作」听起来对称，但若业务里它有方向（总包与分包），对称化会永久抹掉方向信息。

\textbf{反对称性}（Asymmetric）：R(a,b) 成立则 R(b,a) \textbf{必不}成立。
\texttt{directlyPrecedes} 应当反对称：多源数据合并时，
两条方向相反的断言一旦同时出现，推理器立即报不一致——
把静默的数据冲突变成响亮的红灯，这正是声明反对称的价值。
注意辨析：\textbf{反对称}（保证不对称）不同于\textbf{不声明对称}（不保证对称）——
在开放世界下，「没说它对称」永远不等于「它不对称」，想要保证就必须声明。

\textbf{自反性}（Reflexive）：每个个体与自身都有 R 关系。
业务建模中全局自反极少真正需要；经典用例来自部分论（mereology）的约定：
「一切皆为自身之部分」，\texttt{partOf} 自反，而「真部分」\texttt{properPartOf} 则相反。
误用后果：全局自反会给论域里\textbf{每一个}个体加上自环——包括与该属性毫不相干的个体，
徒增噪声与推理负担。若只是想表达「该物料与自身兼容」这类局部事实，
对具体个体直接断言即可，犯不上动用全局公理。

\textbf{非自反性}（Irreflexive）：任何个体与自身都\textbf{没有} R 关系。
\texttt{directlyPrecedes} 是正例：没有工序直接先于它自身。
价值与反对称同理：数据集成时混进一条 \texttt{directlyPrecedes(P5, P5)} 的脏数据，
有非自反公理就是一致性错误，没有它就静默通过——
开放世界下，推理器绝不会「觉得这看起来不对劲」，它的一切判断都必须有公理依据。
另须知道一条 OWL 2 的硬规定：反对称、非自反这类声明只允许加在\textbf{简单属性}上——
传递属性（以及由属性链推出的属性）不属于简单属性，
所以传递的 \texttt{precedes} 不能再声明反对称或非自反，
而不传递的 \texttt{directlyPrecedes} 可以——又一次看到双属性模式的好处。
这条看似琐碎的限制背后是可判定性的边界，2.6 节再谈。

七种特性汇总如下表：

\begin{center}
\begin{tabularx}{\linewidth}{@{}lXXX@{}}
\toprule
\textbf{特性} & \textbf{语义一句话} & \textbf{制造业正例} & \textbf{典型误用后果} \\
\midrule
函数性 & 主体至多一个值 & \texttt{locatedIn}、序列号 &
\texttt{operates} 函数化：多机工人导致设备被推定同一 \\
逆函数性 & 值至多一个主体 & \texttt{hasRFIDTag} &
错误共享值引发错误合并，连锁难回滚 \\
传递性 & 关系沿链传播 & \texttt{precedes}、\texttt{partOf} &
「直接前驱」传递化，「直接」语义尽失 \\
对称性 & 双向自动成立 & \texttt{adjacentTo} &
\texttt{precedes} 对称化，时序约束全面作废 \\
反对称性 & 反向必不成立 & \texttt{directlyPrecedes} &
漏声明时反向脏数据静默通过 \\
自反性 & 个体皆自连 & \texttt{partOf}（部分论约定） &
全局自反带来无谓自环与推理负担 \\
非自反性 & 个体皆不自连 & \texttt{directlyPrecedes} &
漏声明时自环脏数据无法暴露 \\
\bottomrule
\end{tabularx}
\end{center}

收束一句：属性特性不是建模工具里的「可选复选框」，而是\textbf{推理行为的开关}。
每勾选一项之前问两个问题：这条性质在领域里\textbf{恒真}吗
（是「必然如此」，而不是「通常如此」——「通常」是 2.9 节默认推理的领地）？
它将与哪些可能的数据错误相互作用？想清楚再签字。

\section{公理：知识的硬约束}

同样一棵设备分类树，放进 Excel 是「目录」，放进本体才可能成为「理论」——差别全在公理。
知识组织有一条公认的能力阶梯：受控词表（一组规范名词）、
分类法（名词排成层次）、本体（层次之上加公理）。
前两级机器只能「存」，到第三级机器才开始「懂」：
公理让机器第一次有了说「不」的能力——拒绝矛盾的数据、推出没人录入过的结论。
本节按表达力从弱到强过一遍主要公理类型，最后讨论「公理写多强」这个工程权衡。

\subsection{子类公理与等价公理}

\textbf{子类公理}（\(\sqsubseteq\)）给出\textbf{必要条件}：
\texttt{CNCLathe \(\sqsubseteq\) \(\exists\)hasController.CNCController}
读作「凡数控车床必有数控系统」——是数控车床就必须满足；但满足了未必是数控车床。
只用必要条件刻画的类叫\textbf{原始类}（primitive class），
个体归不归它，要靠人工断言，推理器只会向上传播、不会自动归入。

\textbf{等价公理}（\(\equiv\)）给出\textbf{充分必要条件}，得到\textbf{定义类}（defined class）：
推理器可以双向工作——已知个体属于该类，推出它满足全部条件；
已知个体满足全部条件，\textbf{自动}把它归入该类。
后一个方向就是 2.2 节承诺的「自动交叉分类」引擎。看包内教学片段里的例子：

\begin{codebox}
\codeline{\textcolor{cmtgray}{\#\ ==============================}}
\codeline{\textcolor{cmtgray}{\#\ 1.\ 等价类定义\ (Equivalent\ Class)}}
\codeline{\textcolor{cmtgray}{\#\ ==============================}}
\codeline{\textcolor{cmtgray}{\#\ 用于定义类的充分必要条件，明确类的本质特征}}
\codeblank{}
\codeline{\textcolor{cmtgray}{\#\ 定义"合格产品"：}}
\codeline{QualifiedProduct\ \(\equiv\)\ Product\ AND\ (\(\forall\)\ inspection.hasResult(inspection,\ "Pass"))}
\codeline{\textcolor{cmtgray}{\#\ 即：合格产品是"属于Product类"且"所有检验的结果均为通过"的产品}}
\codeblank{}
\end{codebox}
\color{black}

\assetsource{ch02}{description-logic}

「合格产品」被定义为：是产品，且所有检验结果均为通过。
有了这条 \(\equiv\)，一件产品只要满足右边的判据，无须任何人手工打标，
推理器自动将其分类为合格产品；反过来，被断言为合格产品的个体，
推理器也敢替你担保它的每一次检验都通过——这正是「定义」与「描述」的区别。
不过请留心：这个定义里埋着一个相当危险的陷阱——
「所有检验均通过」对一件\textbf{从未检验过}的产品意味着什么？答案在 2.4.4 节揭晓。

\subsection{不相交公理}

\textbf{不相交公理}（disjointness）声明两个类没有共同实例：
\texttt{Equipment \(\sqcap\) Material \(\sqsubseteq\) \(\bot\)}——
\(\bot\) 是空类，整句读作「既是设备又是材料的东西不存在」。

不相交公理的地位远比它的长相重要：\textbf{在开放世界下，它几乎是错误检测的唯一弹药}。
没有它：某个体同时被断言为设备和材料，推理器毫无意见——
也许世上真有东西既是设备又是材料呢？没有公理说不行，那就是「可能」。
有了它：同样的双重断言立即触发不一致。
2.3 节域值域的例子已经演示过这一点：订单被推成产品，
只有「订单与产品不相交」在场时才会报警。

实践经验是：分类树同一层的兄弟类，默认就该考虑两两不相交
（OWL 提供 \texttt{AllDisjoint\allowbreak Classes}，一条声明覆盖一组类）；
但强度要随数据来源的可控程度调节——集成的外部数据源越多、口径越杂，
过强的不相交越容易把「分类口径差异」误判成「逻辑矛盾」，
让一致性检查沦为狼来了。先从最不可能混淆的顶层大类做起，逐层向下收紧，是稳妥的节奏。

一个有真实风味的场景可以说明这笔投资的回报。
集成物料清单与设备台账两套数据时，「冷却液站」在一边被录成辅助设备，
在另一边被挂上了物料编码。若「设备与材料不相交」在位，
合并当晚的一致性检查就会把这对冲突拎出来、精确到涉事个体；
若不在位，两套口径就在图里安静地并行数月，
直到某张报表把冷却液站算进物料库存才东窗事发。
公理的价值不在优雅，而在\textbf{把发现错误的时间从「数月后」提前到「当晚」}。

\subsection{基数约束}

\textbf{基数约束}（cardinality restriction）限定属性值的个数：
恰好 n 个、至少 n 个、至多 n 个。
台账场景的标配是「恰好一个序列号」：\texttt{Equipment \(\sqsubseteq\) =1 hasSerialNumber}。
更有表达力的是\textbf{限定基数}（qualified cardinality）——不光数个数，还限定值的类型：
\texttt{FiveAxisMachine \(\equiv\) CNCMachine \(\sqcap\) \(\ge\)2 hasRotaryAxis.RotaryAxis}，
「五轴机床就是带至少两根旋转轴的数控机床」——数的是「旋转轴这一类的值」，
而不是所有轴。

两条提醒。第一，基数约束与「OWL 不设唯一名假设」相互作用：
「至多一个」遇上两个名字，推理器的第一反应是把两个名字合并成一个个体，
而不是报错（机制与 2.3 节函数性完全相同，系统性讨论见 2.5 节）。
第二，基数约束属于推理代价较高的构造，而且 OWL 2 规定它\textbf{不能套在传递属性上}
（又是「简单属性」限制）——设计属性时就要预留好哪条用于记录、哪条用于传递推理。

\subsection{存在限制、全称限制与「全称陷阱」}

类表达式真正的表达力来自两种受限量词。看片段：

\begin{codebox}
\codeline{\textcolor{cmtgray}{\#\ ==============================}}
\codeline{\textcolor{cmtgray}{\#\ 2.\ 全称限制与存在限制}}
\codeline{\textcolor{cmtgray}{\#\ ==============================}}
\codeblank{}
\codeline{\textcolor{cmtgray}{\#\ 全称限制（Universal\ Restriction，\(\forall\)）：}}
\codeline{\textcolor{cmtgray}{\#\ \(\forall\)R.C\ 表示与该对象存在R关系的所有对象，都属于C类}}
\codeline{\(\forall\)R.C}
\codeblank{}
\codeline{\textcolor{cmtgray}{\#\ 存在限制（Existential\ Restriction，\(\exists\)）：}}
\codeline{\textcolor{cmtgray}{\#\ \(\exists\)R.C\ 表示该对象至少存在一个R关系的对象属于C类}}
\codeline{\(\exists\)R.C}
\codeblank{}
\codeline{\textcolor{cmtgray}{\#\ 使用构造器定义"高端设备"：}}
\codeline{HighEndEquipment\ \(\equiv\)\ Equipment\ \(\sqcap\)\ (price\ \(\ge\)\ 1000000)\ \(\sqcap\)\ (\(\forall\)hasFeature.AdvancedFeature)}
\codeline{\textcolor{cmtgray}{\#\ 即：高端设备是"属于设备类"、"价格大于等于100万"且"所有特征都是高级特征"的设备}}
\codeblank{}
\end{codebox}
\color{black}

\assetsource{ch02}{description-logic}

\textbf{存在限制}（existential restriction）\(\exists\)R.C 刻画「拥有什么」：
至少有一个 R 关系对象属于 C——「数控车床拥有数控系统」。
\textbf{全称限制}（universal restriction）\(\forall\)R.C 刻画「只有什么」：
所有 R 关系对象都属于 C——「高端设备的特征全是高级特征」。
示例中的「高端设备」同时用了三个条件：是设备、价格不低于一百万、所有特征都是高级特征。
看起来无懈可击——直到你想起逻辑课上那个不起眼的角落。

\begin{noteblock}
\textbf{全称陷阱}：\(\forall\)R.C 对「根本没有任何 R 断言」的个体\textbf{自动为真}。
逻辑学称之为空虚真（vacuous truth）：「我的所有法拉利都停在火星上」严格为真——
只要我一辆法拉利都没有。\(\forall\) 不蕴含 \(\exists\)：「所有都是」不保证「至少有一」。
\end{noteblock}

把陷阱代入上面的定义：一台价格 120 万、但 ABox 里\textbf{没有任何}
\texttt{hasFeature} 记录的普通机床，会被推理器堂堂正正地分类为高端设备——
「它的所有特征都是高级特征」：零个特征，全称命题空虚地成立。
再代入「合格产品」：一件从未检验过的产品，「所有检验都通过」为真，
直接被分类为合格产品——在质量管理语境下，这是事故级的推理结论。
雪上加霜的是开放世界：「没有检验记录」并不等于「没检验过」（可能检了没录入），
但全称限制不管这些——它只看图里有什么，空虚满足是实打实的分类结果。

修正的模式是固定的：\textbf{让 \(\exists\) 与 \(\forall\) 成对出现}——

\begin{noteblock}
\centering
QualifiedProduct \(\equiv\) Product \(\sqcap\)
\(\exists\)hasInspection.PassedInspection \(\sqcap\)
\(\forall\)hasInspection.PassedInspection
\end{noteblock}

「至少检过一次，并且每一次都通过」。
从此养成职业习惯：每写一个 \(\forall\)，先自问一句——
\textbf{当关系为空的时候，我希望这个个体算不算？}
答案是「不算」的场合（绝大多数业务定义都是），就补上 \(\exists\)。

\subsection{矛盾、一致性与公理强度的权衡}

公理多了，彼此打架在所难免。包内教学片段给出一对典型矛盾：

\begin{codebox}
\codeline{\textcolor{cmtgray}{\#\ ==============================}}
\codeline{\textcolor{cmtgray}{\#\ 3.\ 公理示例\ (Axiom\ Examples)}}
\codeline{\textcolor{cmtgray}{\#\ ==============================}}
\codeblank{}
\codeline{\textcolor{cmtgray}{\#\ 一致性公理示例（此两条公理相互矛盾）：}}
\codeline{\textcolor{cmtgray}{\#\ 公理1：所有设备都必须位于某个工作单元}}
\codeline{\(\forall\)x:\ Equipment(x)\ \(\rightarrow\)\ locatedIn(x,\ WorkCell)}
\codeline{\textcolor{cmtgray}{\#\ 公理2：存在不位于任何工作单元的设备}}
\codeline{\(\exists\)x:\ Equipment(x)\ \(\land\)\ \(\lnot\)(\(\exists\)w:\ locatedIn(x,\ w))}
\codeblank{}
\end{codebox}

\assetsource{ch02}{description-logic}

公理 1 要求每台设备都位于某个工作单元，公理 2 却断言存在不位于任何工作单元的设备——
二者不可能同时为真。推理器的\textbf{一致性检查}（consistency checking）
会确定性地发现这种冲突并定位涉事公理。这里有一条必须烙进流程的纪律：
\textbf{不一致的知识库里，任何命题都能被「证明」}（逻辑爆炸：矛盾蕴含一切），
所以红灯亮起期间的一切推理输出都不可信——先修复，再使用，没有例外。

最后回答本节开头悬着的问题：公理该写多强？公理是双刃剑：
写得越强，错误暴露越早、可推出的隐含知识越多；
但推理开销越大、对脏数据越敏感、修改维护的门槛越高。三条经验值得记下。
其一，\textbf{能用 \(\sqsubseteq\) 就不用 \(\equiv\)}：
定义类是自动分类的引擎，也是推理负担的大户，只给真正需要自动归类的概念；
普通概念用必要条件描述就够了（包内教学片段 \texttt{description-logic}
第 4 节记录的正是这条优化经验）。
其二，\textbf{不相交先窄后宽}：从顶层大类开始声明，随数据质量的提升逐层加密。
其三，\textbf{每条公理留注释}：公理是领域的「制度条文」，
要写明业务出处与责任人，否则半年后没人敢动它。
至于公理从哪里来——如何从业务问题系统性地提炼公理，那是第 3 章建模方法论的主线；
这些数学记号如何落成机器可读的 OWL 语法，则是第 4 章的任务。

本节符号密集，离开之前收一张速查表，供后续章节随时回查：

\begin{center}
\begin{tabularx}{\linewidth}{@{}llXX@{}}
\toprule
\textbf{公理} & \textbf{记号} & \textbf{读法} & \textbf{工程角色} \\
\midrule
子类 & \texttt{C \(\sqsubseteq\) D} & 凡 C 必是 D（必要条件） &
层次骨架；结论沿层次向上传播 \\
等价 & \texttt{C \(\equiv\) D} & 互为充分必要条件 &
定义类；自动分类的引擎 \\
不相交 & \texttt{C \(\sqcap\) D \(\sqsubseteq\) \(\bot\)} & C 与 D 无共同实例 &
开放世界下错误检测的主要弹药 \\
存在限制 & \(\exists\)R.C & 至少一个 R 关系对象属于 C &
刻画「拥有什么」 \\
全称限制 & \(\forall\)R.C & 所有 R 关系对象都属于 C &
刻画「只有什么」；对空关系自动满足 \\
基数限制 & =1 R 等 & 值的个数恰好、至少、至多 &
键与容量类约束；受简单属性限制 \\
键 & \texttt{HasKey} & 一组属性值联合唯一确定个体 &
跨源数据对齐（第 4 章展开） \\
\bottomrule
\end{tabularx}
\end{center}

\section{实例与断言：ABox 的世界}

TBox 设计得再漂亮，价值也要等 ABox 灌入数据才能兑现——
而 ABox 的麻烦，几乎全和「名字」有关：谁和谁重名？谁和谁其实是同一个？
没提到的名字算不算不存在？本节处理这些看似琐碎、实则左右系统行为的问题。

\subsection{三种断言，外加一种「说不」}

ABox 由断言（assertion）构成，基本款有三种：
\textbf{类断言}——\texttt{CNCLathe(Lathe\_001)}，「这台个体是数控车床」；
\textbf{对象属性断言}——\texttt{producedBy(Product\_A, Lathe\_001)}；
\textbf{数据属性断言}——\texttt{hasPower(Lathe\_001, 15.5)}。
OWL 2 还提供\textbf{否定属性断言}（negative property assertion），
显式声明「\texttt{Lathe\_001} \textbf{不}位于 B 区」。
在封闭世界系统里这是废话——没记录就当不在；
在开放世界下它却是珍贵的显式知识：「未知」与「已知为否」是两种完全不同的状态，
这条区别将贯穿 2.8 节。

\subsection{个体命名：IRI 是身份证}

每个个体由一个 IRI（国际化资源标识符）全局唯一标识。命名风格历来两派：
\textbf{编号型}（含义中立，如 \texttt{EQ-2023-0117}）与
\textbf{语义型}（一眼可读，如 \texttt{No3Lathe}）。工程上推荐编号型，理由只有一个词：
\textbf{不变性}。IRI 是个体的身份证号，要求终身不变；
而「3 号」是车间编号、会重排，「车床」是用途、会改装——
一切可能变化的信息都应该放进属性断言，而不是焊死在标识符里。
人类可读性交给标签解决，与 2.2 节类命名的纪律一致。

\subsection{同名异指与异名同指}

集成两家分厂的数据时，两边都有「3 号车床」——\textbf{同名异指}。
这一半好办：IRI 自带命名空间，
厂 A 与厂 B 各用各的前缀，字符串重名根本不构成冲突。
难办的是另一半：ERP 里的 \texttt{EQ-0117} 与 MES 里的 \texttt{LATHE-3}
其实是车间里同一台机器——\textbf{异名同指}。
两个系统、两套编码、两个 IRI，指着同一个物理对象。
要不要合并？怎么合并？不合并会怎样？这就引出 OWL 实例层最重要的一项语义决策。

\subsection{唯一名假设：OWL 默认不采用}

数据库世界有一条不言自明的铁律：不同主键就是不同实体。
这叫\textbf{唯一名假设}（Unique Name Assumption, UNA）。
\textbf{OWL 默认不采用 UNA}：两个不同的 IRI，在没有进一步信息之前，
既可能指两台设备，也可能是同一台设备的两个名字——推理器对此保持沉默。

为什么这样设计？OWL 生于 Web：知识来自无数互不知情的源头，
同一个事物在不同源头有不同名字是\textbf{常态}而非异常。
若贸然假定「名字不同即事物不同」，任何跨源知识合并都会在第一步就制造海量假矛盾。
放弃 UNA 与开放世界假设是同一套世界观的两面：对没说过的事保持克制。

不采用 UNA 的三个直接后果，前文已埋了两处伏笔，这里集中收割。
\textbf{后果一}：函数性、基数这类「数数」的公理遇到两个名字时，
推理器优先怀疑「这俩是同一个」而不是「数据错了」——
2.3 节工人操作两台设备被推定 \texttt{sameAs} 的机制根源就在这里。
\textbf{后果二}：计数查询变得保守——「厂里有几台设备」，
推理器只能回答「至少有……」，因为它无法排除名字之间的潜在重合。
\textbf{后果三}：想要数据库式的报错行为，必须\textbf{显式声明互异}：
\texttt{owl:differentFrom} 两两声明，或 \texttt{AllDifferent} 一次声明一批。
实战惯例：从台账批量导入个体时，顺手生成一条覆盖全部个体的
\texttt{AllDifferent}——相当于在 OWL 里手工恢复 UNA，
之后函数性与基数公理才会像约束一样工作。

\subsection{sameAs 的威力与代价}

回到异名同指：\texttt{owl:sameAs} 宣布两个 IRI 指同一个体。
声明 \texttt{EQ-0117 owl:sameAs LATHE-3} 之后，两边的断言互相继承：
ERP 侧的维修记录与 MES 侧的产能参数瞬间汇于一身——这是跨系统知识融合最有力的一招。
威力与代价同源：\texttt{sameAs} 的传播没有刹车，
错误的合并会让两台设备的参数互相污染，并顺着推理链扩散到所有下游结论；
在大型知识库中，清理一条错误 \texttt{sameAs} 的代价远高于补一百条漏数据。
因此经验法则是：\textbf{sameAs 宁缺毋滥}——
优先靠逆函数属性、键公理这类「证据驱动」的机制让推理器\textbf{算出}同一性，
人工断言只留给有据可查的少数场合，且每条都要记录来源。

\subsection{类还是实例：型号问题}

ABox 一侧还有一个绕不过的经典难题：\textbf{型号}。「CNC-2000」到底是类还是实例？
对设备厂商，它是产品目录里的一行——一个实例：有发布日期、有定价、有说明书版本；
对使用它的车间，它是 17 台在役机器的共同模板——一个类。两种直觉都对，这正是难点。

常见处理有三种。其一，\textbf{型号建成类}
（\texttt{CNC2000 \(\sqsubseteq\) CNCLathe}），每台机器是它的实例——
当你主要关心「这一型的机器都怎么样」时最顺手。
其二，\textbf{型号建成实例}，机器通过 \texttt{hasModel} 指向它——
当你需要对型号本身大量断言（认证日期、停产时间、适用工艺包）时更自然。
其三，OWL 2 的 punning 机制允许同一个名字同时扮演类与实例两个角色，
代价是两个角色在逻辑上互不相通（第 4 章细说）。
选择的判据可以浓缩成一个问题：\textbf{你更常对「型号本身」说话，
还是对「这一型的每台机器」说话？}前者选实例，后者选类，
两头都重则考虑 punning，或干脆拆成「型号个体 + 机型类」两个建模元素并用属性挂钩。
没有标准答案，只有与查询需求的匹配度——这也预演了第 3 章的一个主题：
建模决策由能力问题（你要回答什么）驱动，而不是由「世界的真相」驱动。

\section{描述逻辑导引：从 ALC 到 SROIQ}

前面五节，我们一直在「用」一套记号：\(\sqsubseteq\)、\(\sqcap\)、\(\equiv\)、
\(\exists\)、\(\forall\)、\(\bot\)……现在正式介绍它们所属的语言家族——
\textbf{描述逻辑}（Description Logic, DL），OWL 的数学地基。
说「导引」，是因为本节不证明定理，只建立三样东西：
一张家族地图、一个三盒结构、一种关于「表达力与代价」的直觉。

\subsection{知识库的三盒结构}

一个 DL 知识库由两到三个「盒子」组成。
\textbf{TBox}（术语盒）装类层面的公理：类层次、定义、不相交——2.2 与 2.4 节的全部内容。
\textbf{ABox}（断言盒）装个体层面的事实——2.5 节的全部内容。
表达力强的家族成员还有第三个盒子：\textbf{RBox}（角色盒，DL 把属性叫「角色」role），
装属性层面的公理——属性层次（\texttt{directlyPrecedes \(\sqsubseteq\) precedes}）、
传递性等七种特性（2.3 节的全部内容），以及\textbf{属性链}
（property chain，例如「设备位于工位，工位部分于车间，则设备位于车间」——
由 \texttt{locatedIn} 与 \texttt{partOf} 链式推出 \texttt{locatedIn}）。
三个盒子分层存放，推理器各有专门算法伺候。

\subsection{基础语言 ALC}

家族的公认基准是 \textbf{ALC}（Attributive Language with Complement）。
它的概念构造器一共五件套：交 \(\sqcap\)、并 \(\sqcup\)、补 \(\lnot\)、
存在限制 \(\exists\)R.C、全称限制 \(\forall\)R.C，外加空类 \(\bot\) 与万有的顶类。
就这五件，已经能写出 2.2 与 2.4 节的大部分骨架：层次、定义类、不相交、存在与全称约束。
写不出的是：基数约束、逆属性、传递性、用个体定义的类。
ALC 的概念可满足性判定是 ExpTime 完全的——最坏情形指数时间，
听起来吓人，但现代推理器靠几十年积累的优化手段，在真实结构的本体上普遍跑得动。
「最坏复杂度」与「实际可用」之间的落差，是 DL 工程的常态，后面还会回到这一点。

顺带把「语义」这个被全章反复使用的词说实。DL 的语义来自模型论：
一个\textbf{解释}（interpretation）做两件事——指定一个论域（所有个体组成的集合），
再把每个类映射为论域的一个子集、把每个属性映射为论域上的一个二元关系。
一条公理在某个解释下「为真」，就是映射结果满足公理的条件
（例如 \(\sqsubseteq\) 两侧映成的子集确实一个包含另一个）；
知识库\textbf{蕴含}某个结论，意思是\textbf{在所有}让全部公理为真的解释中，
该结论\textbf{无一例外}地成立。现在回看 2.2 节那句「在任何可能的状态下」——
它不是修辞，而是「在所有解释中」的通俗说法。
推理器的全部工作，就是替你检查那无穷多个解释的共同点；
可判定性（见下文）则保证这种看似不可能的检查存在有限步算法。

\subsection{字母汤：扩展构造器}

DL 家族的命名像化学分子式：在基准之上，每多一种能力，名字里多一个字母。
常用字母表如下：
\textbf{S} 是 ALC 加传递属性的简写；
\textbf{H} 表示属性层次（Hierarchy）；
\textbf{O} 表示枚举类（nominal——用列举个体的方式造类，如「\{Lathe\_001\}」这种单点类）；
\textbf{I} 表示逆属性（Inverse）；
\textbf{N} 表示非限定基数（只数个数）；
\textbf{Q} 表示限定基数（数「某类的」个数，强于 N）；
\textbf{R} 表示富角色公理——属性链、自反与非自反、不相交属性等；
\textbf{F} 表示函数性；
后缀 \textbf{(D)} 表示支持数据类型（datatype，即数据属性）。
把字母拼起来就是语言名：SHIF(D)、SHOIN(D)、SROIQ(D)——
读者现在可以自行解码：SROIQ 即「ALC + 传递 + 富角色公理 + 枚举类 + 逆属性 + 限定基数」。

\begin{center}
\begin{tabularx}{\linewidth}{@{}lXlX@{}}
\toprule
\textbf{语言} & \textbf{相对前一级的新增能力} & \textbf{推理复杂度} & \textbf{与 OWL 的对应} \\
\midrule
ALC & 布尔构造器、\(\exists\)/\(\forall\) 限制 & ExpTime 完全 &
理论基准，无直接对应 \\
SHIF(D) & 传递、属性层次、逆属性、函数性、数据类型 & ExpTime 完全 &
OWL Lite（OWL 1） \\
SHOIN(D) & 枚举类、非限定基数 & NExpTime 完全 &
OWL DL（OWL 1） \\
SROIQ(D) & 属性链、限定基数、（非）自反、不相交属性 & N2ExpTime 完全 &
OWL 2 DL \\
EL、DL-Lite 等子集 & 反向裁剪：砍掉昂贵构造器 & 多项式时间 &
OWL 2 EL/QL/RL Profile（第 4 章） \\
\bottomrule
\end{tabularx}
\end{center}

这张表要竖着读出两个趋势。
向下，表达力上涨，复杂度从指数、非确定指数一路爬到双重非确定指数——
\textbf{每一格表达力都明码标价}。
最底一行则是反向操作：EL、DL-Lite 这些「轻量级」成员故意砍掉昂贵构造器，
把推理压进多项式时间，专门伺候千万级实例的大规模场景——
它们正是 OWL 2 三个 Profile 的理论原型，选型策略留给第 4 章。

工程师该怎么读这张表？三条阅读法。
第一，复杂度衡量的是\textbf{最坏情形}且针对\textbf{整门语言}：
你的本体只要不踩满全部构造器组合，实际推理时间往往与理论上界相去甚远——
被「双重指数」吓得不敢用 OWL 2 DL，和无视复杂度随手堆公理，是对称的两种错误。
第二，真正的开销大户是少数模式：大量定义类（\(\equiv\)）、
逆属性与基数的密集组合、巨型不相交簇——2.4 节「能用 \(\sqsubseteq\) 不用 \(\equiv\)」
的优化经验，在这张表里有了理论注脚。
第三，当 ABox 膨胀到千万级，与其在 SROIQ 里苦苦调优，
不如检查表达力需求能否塞进某个 Profile——降一级语言，常常换来几个数量级的吞吐。

\subsection{可判定性：为什么是设计底线}

一个逻辑是\textbf{可判定的}（decidable），指存在算法对任何合法输入
都能在\textbf{有限步内}给出确定的是否答案。一阶逻辑不可判定（丘奇与图灵在
1936 年分别证明）：不存在通用算法判定任意一阶公式是否有效——
换句话说，基于全量一阶逻辑的推理器，理论上就无法承诺「一定算得完」。

对生产系统而言这是底线问题：一致性检查若可能永不返回，
任何上线承诺、任何 SLA 都无从谈起。描述逻辑全家族的设计哲学因此可以浓缩成一句话：
\textbf{在保住可判定性的前提下，把表达力推到最远}。
SROIQ 就是这场推到极限的游戏在 OWL 2 时代的成果——再往前迈半步
（比如取消属性链上的限制），可判定性立刻失守。
现在可以回收 2.3 节那条「琐碎」规定了：传递属性不能声明反对称、不能参与基数约束，
不是标准委员会小气，而是\textbf{可判定性的边界恰好画在那里}——
「简单属性」限制是 SROIQ 不掉下悬崖的护栏。

\subsection{推理器拿这些公理做什么}

有了严格语义，推理器能提供一组\textbf{标准推理服务}，名字值得现在就记住：

\begin{itemize}
\item \textbf{一致性检查}（consistency）：整个知识库有没有矛盾——2.4 节的红灯；
\item \textbf{分类}（classification）：算出所有隐含的 \(\sqsubseteq\) 关系，
      把每个类放到层次中的正确位置——2.2 节「让机器算交叉分类」的引擎；
\item \textbf{概念可满足性}（satisfiability）：某个类是否注定为空——
      注定为空的类几乎总是公理冲突的症状；
\item \textbf{实例检查与实现}（instance checking / realization）：
      某个体属于哪些类、最具体的类是哪个——ABox 一侧的分类。
\end{itemize}

这些服务背后的算法（tableau 等）与工程用法，第 5 章推理机制专章展开；
本章只需要建立这样一个分工印象：\textbf{建模者负责把知识写成公理，
推理器负责把公理的全部逻辑后果挖出来}——前者是人的判断，后者是机器的体力。

\section{一阶逻辑与规则：DL 之外的表达力}

描述逻辑写不出「容量更大的故障设备优先维修」。
这不是它的耻辱，而是它为可判定性支付的账单——上一节刚刚算过这笔账。
本节走到墙外看看：一阶逻辑有多大的天地，DL 在其中处于什么位置，
以及工程上如何安全地借用墙外的力量。

\subsection{从 DL 到 FOL：翻译与定位}

\textbf{一阶逻辑}（First-Order Logic, FOL）的构件是：谓词、变量、
全称与存在量词（\(\forall\)、\(\exists\)）、以及与或非等联结词。
DL 的每个构造都能机械地翻译成 FOL：
\texttt{C \(\sqsubseteq\) D} 即「对一切 x：若 C(x) 则 D(x)」；
\(\exists\)R.C 即「存在 y：R(x,y) 且 C(y)」。
所以一个准确的定位是：\textbf{DL 是 FOL 的一个可判定子集}——
大体上相当于把变量个数与量词使用模式严格受限之后剩下的那部分 FOL
（理论上对应所谓二变量片段与守卫片段的交叠地带）。
变量少、量词受限，正是可判定性的来源；
反过来，FOL 里那些 DL 写不出的句子，几乎都在自由地使用多个变量。

\subsection{工程规则的 FOL 表达}

什么样的业务规则会用满 FOL 的自由度？看包内教学片段里的四条：

\begin{codebox}
\codeline{\textcolor{cmtgray}{\#\ ==============================}}
\codeline{\textcolor{cmtgray}{\#\ 工程约束的一阶逻辑表达}}
\codeline{\textcolor{cmtgray}{\#\ ==============================}}
\codeblank{}
\codeline{\textcolor{cmtgray}{\#\ 规则：所有设备都需要定期维护}}
\codeline{\(\forall\)x:\ Equipment(x)\ \(\rightarrow\)\ \(\exists\)m:\ Maintenance(m)\ \(\land\)\ requires(x,\ m)\ \(\land\)\ isPeriodic(m)}
\codeblank{}
\codeline{\textcolor{cmtgray}{\#\ 规则：容量大的故障设备优先维修}}
\codeline{\(\forall\)a,b:\ (Equipment(a)\ \(\land\)\ Equipment(b)\ \(\land\)\ capacity(a)\ >\ capacity(b)\ \(\land\)\ hasFault(a))}
\codeline{\hspace*{4\ttcw}\(\rightarrow\)\ priority(repair(a))\ >\ priority(repair(b))}
\codeblank{}
\codeline{\textcolor{cmtgray}{\#\ 规则：同一时刻设备只能执行一个任务}}
\codeline{\(\forall\)\ e,\ t,\ task1,\ task2:\ (executes(e,\ task1,\ t)\ AND\ executes(e,\ task2,\ t))}
\codeline{\hspace*{4\ttcw}\(\rightarrow\)\ (task1\ =\ task2)}
\codeblank{}
\codeline{\textcolor{cmtgray}{\#\ 规则：热处理必须在焊接之后进行}}
\codeline{\(\forall\)\ p1,\ p2:\ (p1.type\ =\ "HeatTreatment"\ AND\ p2.type\ =\ "Welding")}
\codeline{\hspace*{4\ttcw}\(\rightarrow\)\ precedes(p2,\ p1)}
\codeblank{}
\end{codebox}
\color{black}

\assetsource{ch02}{first-order-logic}

逐条体检，可以精确感受 DL 边界的位置。
\textbf{规则一}（设备需有定期维护）其实没出圈：
它就是 \texttt{Equipment \(\sqsubseteq\) \(\exists\)requires.PeriodicMaintenance}，DL 完全胜任。
\textbf{规则三}（同一时刻只执行一个任务）麻烦在「时刻 t」：
设备、任务、时间构成\textbf{三元关系}，而 DL 的属性天生二元——
要么把「执行」具体化（reification）成一个带三条边的小对象，要么交给规则语言。
\textbf{规则二}（容量大的故障设备优先维修）与\textbf{规则四}（热处理必须在焊接之后）
是真正的圈外人：它们在多个个体之间\textbf{自由共享变量}、还要\textbf{比较数据值}
（容量大小、工序类型）——这正是 FOL 的强项，也是它不可判定的根源所在。
一句直觉：\textbf{DL 擅长「定义一类东西」，FOL 擅长「描述几件东西之间的纠葛」}。

\subsection{Horn 子句：受限而高效的规则}

全量 FOL 不可判定，但工程并不需要全量。
最重要的受限片段是 \textbf{Horn 子句}（Horn clause）：
形如「条件一 且 条件二 且 …… 则 结论」、且结论只许一个原子的规则。
它的直觉形象就是程序员熟悉的 IF--THEN：
没有「或者型结论」（不许说「故障原因是 A 或 B」），没有否定，
推理可以沿着规则链高效地正向传播。
数据库领域的 Datalog 语言就建在 Horn 子句上；
本体世界的对应物是第 5 章的主角 \textbf{SWRL}（OWL 加 Horn 规则）。

感受一条 Horn 规则的样子（用可读的伪写法）：
「若设备 e 位于单元 c，且工人 w 被指派到单元 c，则 w 可操作 e」。
三个变量横跨三个个体，前件两个条件，结论恰好一个原子——
DL 写不出它（变量无法这样自由共享），全量 FOL 又嫌它平凡，Horn 恰到好处。
真实工程规则库的主体，正是成百上千条这种「平凡」规则的集合。
SWRL 若放任变量绑定到任意推理出的匿名个体，会重新跌入不可判定；
工程上通行 \textbf{DL-safe} 限制——规则变量只绑定到 ABox 里有名有姓的个体——
以一点表达力换回可判定性，延续的仍是 2.6 节那条设计哲学。

于是全章反复出现的分工格局可以正式立牌：
\textbf{骨架交给 DL}（类层次、定义、约束——可判定、自动分类、全局一致性），
\textbf{肌肉交给规则}（跨个体关联、数值比较、三元以上关系——局部、高效、贴业务）。
两者如何在同一个系统里协同，第 5 章给出完整方案。

离开本节前回答一个常见质疑：「这些规则我用 Python 半小时也能写完，为什么要走逻辑？」
区别不在能不能实现，而在知识以什么形态存在。写成逻辑规则，它是\textbf{声明式}的知识：
可被审查（领域专家读得懂前件与后件，不需要会编程）、
可被组合（多条规则自动协同，无须人工编排执行顺序）、
可被解释（每个结论都能回溯到触发它的规则与事实）。
写成代码，它是过程：知识重新隐没回实现细节，
评审、复用、解释都要透过代码考古完成——恰恰退回了第 1 章描述的
「知识锁在系统与脑子里」的起点。把业务规则保持在声明层，是本体工程从头到尾的立场。

\section{开放世界与封闭世界}

周五下午，你写了个查询找「下周没有排任务的设备」，准备安排保养窗口。
查询一跑，结果为空。是全厂设备真的都排满了，还是查询写错了？
都不是。你撞上的是 OWL 最著名的「性格」——开放世界假设。
这一节把第 1 章埋下的那条线索（「无记录」不等于「不能」）正式收回来。

\subsection{两种世界观}

\textbf{开放世界假设}（Open World Assumption, OWA）：
知识库未记载的命题，真值\textbf{未知}——既不当真，也不当假。OWL 与描述逻辑采用。
\textbf{封闭世界假设}（Closed World Assumption, CWA）：
未记载即为\textbf{假}——查无此记录，就是没有；
其实现机制叫\textbf{否定即失败}（Negation as Failure, NAF）。
关系数据库与 Prolog 采用。看片段里的最小对照：

\begin{codebox}
\codeline{\textcolor{cmtgray}{\#\ ==============================}}
\codeline{\textcolor{cmtgray}{\#\ 3.\ 开放世界假设与封闭世界假设}}
\codeline{\textcolor{cmtgray}{\#\ ==============================}}
\codeline{\textcolor{cmtgray}{\#\ OWA\ (Open\ World\ Assumption)\ \ ——\ OWL\ /\ 描述逻辑采用}}
\codeline{\textcolor{cmtgray}{\#\ \ \ 未声明的事实\ =\ 未知（可能为真，也可能为假）}}
\codeline{\textcolor{cmtgray}{\#\ CWA\ (Closed\ World\ Assumption)\ ——\ 关系数据库\ /\ Prolog\ 采用}}
\codeline{\textcolor{cmtgray}{\#\ \ \ 未声明的事实\ =\ 假（否定即失败，Negation\ as\ Failure）}}
\codeblank{}
\codeline{\textcolor{cmtgray}{\#\ 示例：知识库中只有一条事实}}
\codeline{producedBy(Product\_A,\ Lathe\_001)}
\codeblank{}
\codeline{\textcolor{cmtgray}{\#\ 问题：Product\_A\ 是否由\ Mill\_002\ 生产？}}
\codeline{\textcolor{cmtgray}{\#\ OWA\ 回答：未知（除非显式声明\ \(\lnot\)producedBy(Product\_A,\ Mill\_002)）}}
\codeline{\textcolor{cmtgray}{\#\ CWA\ 回答：否（知识库未记录即视为假）}}
\codeblank{}
\codeline{\textcolor{cmtgray}{\#\ 工程影响：}}
\codeline{\textcolor{cmtgray}{\#\ 在\ OWL\ 中校验"设备未分配任何任务"这类否定条件时，}}
\codeline{\textcolor{cmtgray}{\#\ 需借助\ SPARQL\ 的\ FILTER\ NOT\ EXISTS（局部封闭世界）实现}}
\codeblank{}
\end{codebox}
\color{black}

\assetsource{ch02}{reasoning-examples}

同一条知识、同一个问题，两种世界观给出不同回答。
OWA 的道理一句话能讲完：\textbf{知识库永远不完备}——
没说 \texttt{Product\_A} 不是 \texttt{Mill\_002} 生产的，
也许只是那条记录还没进来；缺一条边，不能当作反面证据。
Web 规模的知识集成必须这样保守，否则任何两个数据源一合并就互相「证伪」。
而数据库管理的是边界明确的业务闭环——订单表里没有这张订单，那就是没有——CWA 同样合理。
\textbf{两种假设没有对错，只有适用域}；危险出在混用而不自知。

有读者会想起，数据库其实也有「第三种状态」：NULL，以及围绕它的三值逻辑。
形似而神异。NULL 是\textbf{字段级}的例外标记，语义出了名地含混
（是未知？不适用？还是没来得及填？），聚合与连接时的行为更是面试题常客；
OWA 的「未知」则是\textbf{系统级}的默认立场——不是某个字段特殊，
而是整个知识库对一切未记载的命题统一保持克制。
一句话：数据库把「未知」当例外补丁，OWL 把「未知」当世界的常态。
理解这层差别，跨系统对接时就不会把 NULL 与「图中无此三元组」想当然地互译——
前者往往被各业务系统赋予了私有含义，后者在语义上什么都没说。

\subsection{四象限真值对照}

把一个命题 P 在知识库中的证据状态分成四种情形，两种世界观的差异一目了然：

\begin{center}
\begin{tabularx}{\linewidth}{@{}XccX@{}}
\toprule
\textbf{知识库的证据状态} & \textbf{OWA 判定} & \textbf{CWA 判定} & \textbf{工程含义} \\
\midrule
能证明 P（有断言或可推出） & 真 & 真 & 两种世界观一致 \\
能证明 P 不成立（显式否定或不相交推出） & 假 & 假 &
一致；但 CWA 系统极少显式存储否定 \\
P 与其否定都证不出（无记录） & \textbf{未知} & \textbf{假} &
全部分歧所在：「无记录」的解释权 \\
P 与其否定都能证出 & 不一致，全库报警 & 视实现而定，常被静默吞掉 &
OWA 把矛盾当一级错误对待 \\
\bottomrule
\end{tabularx}
\end{center}

真正的分水岭只有第三行：\textbf{对「无记录」的解释}。
CWA 把它折叠成「假」，二值世界，痛快但武断；
OWA 让它保持为第三种状态「未知」，三值世界，保守但诚实。

\subsection{「无记录不等于否」：对 AI 产品的意义}

回看第 1 章那个 AI 对照例子：本体校验路径的回答是
「知识库中\textbf{没有}该设备可加工钛合金\textbf{的记录}」，
而不是「该设备\textbf{不能}加工钛合金」。现在读者拥有了完整的理论装备：
前者报告的是第三行的「未知」，后者僭越成了第二行的「假」——一词之差，语义级别完全不同。

对构建 AI 产品的含义是直接的：基于本体回答问题的系统，
应当把三值状态（证实、证伪、无记录）\textbf{透传}到产品层，
而不是在中间某层悄悄折叠成二值。
「未知」触发的正确动作是：检索更多数据源、向用户追问、或如实告知边界——
唯独不是替用户编一个「否」（更不是编一个「是」）。
大语言模型的幻觉问题之所以与本体互补，正在于此：
LLM 的统计本能是「永远给出一个流利的答案」，
本体的逻辑本能是「分得清知道与不知道」。
落到产品对话上，差别一目了然。用户问设备能力助手：「7 号铣床能加工钛合金吗？」
合格的三种回答是：「能——能力清单第 12 条，认证日期某年某月」（证实）；
「不能——其主轴功率低于钛合金工艺下限，依据工艺约束第 3 条」（证伪）；
「知识库中没有该设备与钛合金工艺的任何记录，已转人工确认」（未知）。
不合格的回答只有一种，却最常见：把第三种情形包装成语气自信的前两种之一。
把这条原则落成 LLM 加本体的工程架构，是第 8 章的主题。

\subsection{局部封闭世界：SPARQL 实战}

OWA 是本体语义层的假设，但它并不禁止你在\textbf{查询层}对特定数据「局部封闭」
（Local Closed World）。依据很朴素：当你\textbf{确知}某类数据是完备的——
本厂设备台账齐全、排产系统的任务分配是唯一事实源——
就可以理直气壮地把「查无记录」当「没有」。SPARQL 的惯用写法：

\begin{noteblock}
\texttt{SELECT ?e WHERE \{}\\
\texttt{\ \ ?e rdf:type :Equipment .}\\
\texttt{\ \ FILTER NOT EXISTS \{ ?t :assignedTo ?e \}}\\
\texttt{\}}
\end{noteblock}

这条查询找出「图中不存在任务指派记录」的全部设备——
开头那个保养窗口问题的正解。\texttt{FILTER NOT EXISTS} 检查的是
「图里匹配不到这个模式」，本质是查询层的否定即失败；
它不改变本体的 OWA 语义，只是在这一次查询里、对这一个谓词，临时拉上了世界的边界。

两条纪律让这个利器不伤手。
其一，\textbf{封闭的前提是完备}：对一个采集不全的谓词做 \texttt{NOT EXISTS}，
得到的是系统性误导——「没录入维护记录」会被当成「从未维护」。
其二，\textbf{封闭范围要文档化}：哪些谓词、在哪些查询里、依据什么完备性承诺
按封闭世界解释，应当写进系统的语义边界说明——口头约定迟早出事。
顺带一条同族提醒：把推理结果物化进关系数据库之后，
下游消费者会不自觉地用 CWA 读它（数据库嘛），这条语义换轨同样要写清楚。

\section{单调与非单调推理}

「设备默认可用，除非它故障了。」——这句车间常识，OWL 居然写不出来。
为什么写不出？写不出的能力值多少钱？工程上又该怎么办？这是本章最后一组问题。

\subsection{单调性：定义与价值}

一个推理体系是\textbf{单调的}（monotonic），如果知识只增不减时，结论也只增不减：
向知识库添加任何新公理、新事实，旧结论一条不少地保持成立。
形式地说：知识库变大，其逻辑后果集合不会变小。
经典逻辑——一阶逻辑、描述逻辑——都是单调的；OWL 推理器自然也是。看片段：

\begin{codebox}
\codeline{\textcolor{cmtgray}{\#\ ==============================}}
\codeline{\textcolor{cmtgray}{\#\ 1.\ 单调推理\ (Monotonic\ Reasoning)}}
\codeline{\textcolor{cmtgray}{\#\ ==============================}}
\codeline{\textcolor{cmtgray}{\#\ 定义：新增知识只会增加结论，不会推翻已有结论}}
\codeline{\textcolor{cmtgray}{\#\ 经典逻辑（一阶逻辑、描述逻辑）都是单调的}}
\codeline{\textcolor{cmtgray}{\#\ OWL\ DL\ 推理器（Pellet\ /\ HermiT）遵循单调性假设}}
\codeblank{}
\codeline{\textcolor{cmtgray}{\#\ 已知公理（TBox）：}}
\codeline{CNCLathe\ \(\sqsubseteq\)\ CNCMachine\ \ \ \ \ \ \ \ \ \ \ \ \ \#\ 数控车床是数控机床}
\codeline{CNCMachine\ \(\sqsubseteq\)\ ProcessingEquipment\ \ \#\ 数控机床是加工设备}
\codeblank{}
\codeline{\textcolor{cmtgray}{\#\ 已知事实（ABox）：}}
\codeline{CNCLathe(Lathe\_001)\ \ \ \ \ \ \ \ \ \ \ \ \ \ \ \#\ Lathe\_001\ 是一台数控车床}
\codeblank{}
\codeline{\textcolor{cmtgray}{\#\ 单调推导（子类关系传递）：}}
\codeline{CNCMachine(Lathe\_001)\ \ \ \ \ \ \ \ \ \ \ \ \ \#\ 推论1：Lathe\_001\ 是数控机床}
\codeline{ProcessingEquipment(Lathe\_001)\ \ \ \ \#\ 推论2：Lathe\_001\ 是加工设备}
\codeblank{}
\codeline{\textcolor{cmtgray}{\#\ 单调性保证：之后无论向知识库添加什么新事实}}
\codeline{\textcolor{cmtgray}{\#\ （设备故障、位置变更、新的分类公理），}}
\codeline{\textcolor{cmtgray}{\#\ 推论1、推论2\ 永远成立，不会被撤销}}
\codeblank{}
\end{codebox}

\assetsource{ch02}{reasoning-examples}

两条 TBox 公理加一条 ABox 事实，沿子类链推出两条分类结论；
单调性承诺：今后无论发生什么——设备故障、位置变更、本体扩充——这两条结论永不收回。

别小看这个承诺，它是三项工程美德的来源。
\textbf{可审计}：每个结论都有一条不会被未来推翻的证明链，
监管与合规场景（「系统当时为什么这样判」）必需。
\textbf{可缓存、可增量}：已算出的结论永不失效，
推理结果可以物化落库、可以分批计算、可以只对新增数据做增量推理。
一个广泛采用的部署节奏是：夜间批作业对全库做一次完整分类并把结论物化，
白天的查询直接命中物化结果、零推理开销，新到数据只做小步增量，次日夜间归并。
这套节奏的逻辑前提正是单调性——若白天新增的事实可能推翻昨晚的结论，
物化表就成了定时炸弹。
\textbf{可合并}：两个独立演化的知识库合并后，各自原有的结论仍然全部成立——
这是 Web 知识能够分布式生长的逻辑前提，与 OWA、不设 UNA 同属一套世界观。

\subsection{默认与例外：单调世界的盲区}

但工程知识偏偏充满「默认加例外」的形态：设备默认可用，除非故障；
工序默认按标准工时，除非加急；供应商默认合格，除非进了黑名单。
处理这种知识需要\textbf{非单调推理}（non-monotonic reasoning）——
允许新知识推翻旧结论。经典形式化是 Reiter 的\textbf{默认逻辑}（default logic），
看它的最小工作示例：

\begin{codebox}
\codeline{\textcolor{cmtgray}{\#\ 默认规则（Reiter\ 默认逻辑）：}}
\codeline{Equipment(x)\ :\ available(x)\ /\ available(x)}
\codeline{\textcolor{cmtgray}{\#\ 规则解读：如果\ x\ 是设备，且"x\ 可用"与知识库一致（无矛盾信息），}}
\codeline{\textcolor{cmtgray}{\#\ 则默认推出\ x\ 可用}}
\codeblank{}
\codeline{\textcolor{cmtgray}{\#\ 阶段1\ ——\ 知识库仅有：}}
\codeline{Equipment(Lathe\_001)}
\codeline{\textcolor{cmtgray}{\#\ 默认推论：available(Lathe\_001)\ 成立\ \(\checkmark\)}}
\codeblank{}
\codeline{\textcolor{cmtgray}{\#\ 阶段2\ ——\ 新增事实与规则：}}
\codeline{hasFault(Lathe\_001)\ \ \ \ \ \ \ \ \ \ \ \ \ \ \ \#\ Lathe\_001\ 出现故障}
\codeline{\(\forall\)x:\ hasFault(x)\ \(\rightarrow\)\ \(\lnot\)available(x)\ \ \ \#\ 故障设备不可用}
\codeline{\textcolor{cmtgray}{\#\ 撤销推论：available(Lathe\_001)\ 被推翻\ \(\times\)}}
\codeline{\textcolor{cmtgray}{\#\ 对于没有故障记录的其他设备，仍默认其可用}}
\codeline{\textcolor{cmtgray}{\#\ 这种"结论收回"是单调逻辑无法表达的}}
\end{codebox}

\assetsource{ch02}{reasoning-examples}

默认规则读作三段式：\textbf{前提}（x 是设备）、
\textbf{一致性检验}（「x 可用」与现有知识不矛盾）、\textbf{默认结论}（推出 x 可用）。
阶段 1 里知识库只知道 \texttt{Lathe\_001} 是设备，一致性检验通过，默认结论成立；
阶段 2 新增故障事实后，检验失败，先前的结论被\textbf{收回}——
术语叫\textbf{可撤销推理}（defeasible reasoning）。
「结论随知识增长而收回」恰是单调性的反面，
在经典 DL 框架内\textbf{不可表达}——不是推理器不够聪明，是语义层面就被禁止。

非单调逻辑是一个理论深厚的家族（默认逻辑、环境逻辑、回答集编程……），
但工程主流并不直接采用：非单调语义计算代价高，
且「结论取决于你在哪个时刻问」的不稳定性，与本体作为「权威知识源」的定位相冲突。
工程的智慧是绕行而不是硬闯。

绕行之前，先区分两种容易混淆的「撤销」。一种是\textbf{世界变了}：
设备昨天可用、今天故障——事实本身随时间改变。
这是\textbf{数据更新}问题，不需要任何非单调逻辑，
把状态建成显式数据、随事件更新即可（下文路径一）。
另一种是\textbf{当初就信错了}：知识库一直记着这台设备在 A 区，
后来发现是录入错误——这是\textbf{信念修正}（belief revision）问题，
工程上表现为修正数据再重算下游结论（下文路径四）。
真正需要非单调\textbf{逻辑}的，只剩「在不完全信息下先按默认行事」的推理形态；
而工厂里大多数所谓「默认加例外」，仔细一看其实属于前两种——
这正是接下来四条工程路径足以覆盖绝大多数需求的原因。

\subsection{工程实现途径}

实践中消化「默认加例外」有四条成熟路径，按推荐顺序排列：

\begin{itemize}
\item \textbf{状态显式化}（首选）：把「可用」从默认推论降格为显式数据——
      \texttt{hasStatus} 属性由设备管理系统维护，故障即更新，查询按状态过滤。
      用\textbf{数据更新}代替\textbf{信念撤销}，完全留在单调框架内，成本最低；
\item \textbf{查询层局部封闭}：用 2.8 节的 \texttt{FILTER NOT EXISTS}
      把默认规则改写成查询——「没有故障记录的设备」就是「默认可用的设备」。
      默认语义被限制在单次查询的边界内，本体语义不受扰动；
\item \textbf{应用层规则引擎}：支持否定即失败的规则系统
      （Datalog 引擎、生产式规则引擎、SHACL 规则等）跑在本体\textbf{之上}，
      专职处理默认与例外；本体层只存事实与硬公理，两层各守各的语义；
\item \textbf{重算而非撤销}：知识变更后丢弃旧的物化结果、全量或增量重推。
      单调推理器照常工作，「撤销」发生在物化管理层而非逻辑层——
      以计算换语义简单，批处理场景常用。
\end{itemize}

四条路径共享同一个思想，值得装裱起来：
\textbf{把非单调留给应用层与数据层，让本体层保持单调}——
结论稳定、审计简单、合并安全。
第 5 章讨论 SWRL 与规则引擎时，将看到这套分层在真实技术栈中的形态。

\section{本章小结、思考题与文件指引}

\subsection{本章小结}

本章在四块积木与两套逻辑之间往返，核心论点其实只有一个：
\textbf{本体的每一个设计选择，都是写给推理器的承诺}——
层次怎么挂、特性勾不勾、公理写多强、名字合不合并，
机器都会一丝不苟地兑现，包括你不曾预料的那部分。逐条要点如下：

\begin{itemize}
\item 本体 = 词汇（类与属性）+ 约束（公理）+ 事实（实例断言）；
      TBox 与 ABox 两层分离，骨架与血肉独立演化；
\item 类是内涵不是清单；\(\sqsubseteq\) 是任何状态下的集合包含，继承是逻辑推论而非数据拷贝；
      主干单继承，交叉分类交给定义类与推理器；
\item 实例、部分、角色、维度组合都不是子类——第 3 章 OntoClean 把这些直觉升级为检查表；
\item 域与值域是分类公理而非校验规则；想让错误暴露要靠不相交公理，
      想做输入校验要靠约束语言；
\item 七种属性特性是与推理器的合同：误用的代价从「脏数据静默通过」到「个体被错误合并」；
\item \(\forall\) 对空关系自动为真（全称陷阱）：业务定义几乎总要 \(\exists\) 与 \(\forall\) 成对；
\item OWL 不设唯一名假设：函数性与基数遇到多个名字先怀疑同一性；
      \texttt{AllDifferent} 恢复约束行为，\texttt{sameAs} 宁缺毋滥；
\item DL 是 FOL 的可判定子集；字母即构造器，OWL 2 DL 对应 SROIQ(D)；
      表达力与复杂度同涨，可判定性是设计底线；
\item 开放世界：无记录 = 未知而非假；三值状态应透传到产品层；
      完备数据上可用 \texttt{FILTER NOT EXISTS} 局部封闭；
\item OWL 推理单调；「默认加例外」用状态显式化、查询封闭、规则引擎或重算消化，
      本体层守住单调。
\end{itemize}

\subsection{思考题}

以下八道题覆盖本章全部主干概念。建议动笔写出公理与查询，而不只是口头作答；
多数题目的线索散落在正文对应小节与 Semantica 教学资产中，
个别题目需要把两个小节的结论拼起来用——这是有意为之。

\begin{noteblock}
\begin{enumerate}
\item \textbf{层次设计}：有人提议把「备用设备」「在修设备」建成
      \texttt{Equipment} 的两个子类。用外延与内涵、以及「必然是」判据分析问题所在，
      给出替代建模方案，并预测第 3 章 OntoClean 的「刚性」检查会如何评价这个提议。
\item \textbf{域与值域}：已知 \texttt{producedBy} 的域是 \texttt{Product}，
      且本体声明了 \texttt{Order \(\sqcap\) Product \(\sqsubseteq\) \(\bot\)}。
      断言 \texttt{producedBy(Order\_77, Robot\_R2)} 后，推理器会得出什么？
      删去不相交公理后又会怎样？这说明域公理的本质是什么？
\item \textbf{特性配置}：为 \texttt{directlyPrecedes}、\texttt{precedes}、
      \texttt{locatedIn}、\texttt{hasRFIDTag} 四个属性分别选择应声明的特性
      （函数、逆函数、传递、对称、反对称、自反、非自反），
      并指出哪一种「特性组合」会触犯 OWL 2 的简单属性限制。
\item \textbf{全称陷阱}：按定义
      \texttt{QualifiedProduct \(\equiv\) Product \(\sqcap\)
      \(\forall\)hasInspection.PassedInspection}，
      一件从未检验过的产品会被如何分类？为什么？
      写出修正后的定义，并解释 \(\exists\) 与 \(\forall\) 为何要成对出现。
\item \textbf{唯一名假设}：\texttt{locatedIn} 已声明函数性，ABox 中同时存在
      \texttt{locatedIn(EQ1, CellA)} 与 \texttt{locatedIn(EQ1, CellB)}。
      推理器会报错吗？若不会，它会得出什么结论？
      需要补充什么声明，才能让这类数据错误如预期暴露？
\item \textbf{表达力边界}：把「容量更大的故障设备应获得更高维修优先级」写成一阶逻辑；
      指出其中哪些成分超出了 SROIQ 的表达力，并说明第 5 章的规则机制将如何接手。
\item \textbf{局部封闭}：写出一个 SPARQL 查询，找出「没有任何维护记录的设备」；
      说明在什么数据完备性条件下其结果可信，违反该条件时会产生哪一类错误结论。
\item \textbf{非单调落地}：「设备默认可用，故障即不可用，修复后恢复可用」——
      给出两种保持本体层单调性的工程实现方案，
      并从审计难度与实时性两个角度比较它们的优劣。
\end{enumerate}
\end{noteblock}

\subsection{本章包资产与场景指引}

下表资产位于 \texttt{semantica.chapter\_packages.vol1.ch02}，不在书目录中形成副本。
其中四份文本解释理论；真正可复算的是另行绑定的 OWA/CWA 查询、事实、规则与场景。

\begin{center}
\begin{tabularx}{\linewidth}{@{}>{\ttfamily\footnotesize}p{0.36\linewidth}X@{}}
\toprule
\normalfont\textbf{Semantica asset} & \textbf{内容与阅读建议} \\
\midrule
core-concepts & 四块积木的完整示例：类层次、数据与对象属性、上层本体对应；
配合 2.1 至 2.3 节阅读，建议把每条属性读成「一个可问的问题」 \\
description-logic & 等价类、\(\exists\)/\(\forall\) 限制、一对矛盾公理；
配合 2.4 与 2.6 节阅读；末节的概率标注展示了经典 DL 之外的扩展方向 \\
first-order-logic & 工程规则的 FOL 写法与默认规则最小示例；
配合 2.7 与 2.9 节，逐条判断「哪条 DL 写得出」是极好的练习 \\
reasoning-examples & 单调推导、撤销推理、OWA 与 CWA 对照及工程影响；
配合 2.8 与 2.9 节，注意第 4 节对两类推理适用场景的划分 \\
OE-V1-CH02-SCN-REASONING-MODES-001 & 受限正向链与 OWA/CWA 精确 oracle；
运行成功仍不解除完整 DL、默认逻辑与 release 收据缺口 \\
\bottomrule
\end{tabularx}
\end{center}

理论地基至此打完。下一章把镜头从「积木是什么」转向「怎么搭」——
能力问题、七步法与 OntoClean 方法论将把本章的判断直觉变成可验证的工程流程。
在那之前，建议抽半小时通过 package asset 读取上表四份材料：
本章引用的只是节选，资产里那些未被引用的小节与注释，
正是为已经读完正文的你准备的第二遍材料。
